% =============================================================================
\chapter{Conclusion}
\label{ch:conclusion}
% =============================================================================

\section{Summary of Achievements}
\label{sec:conc:ach}

This project successfully designed, trained, evaluated, and deployed
\emph{Eshara}, a real-time bilingual \gls{asl}/\gls{arsl} recognition
system covering letter-level and word-level signs. Four deep models were
trained: two \gls{mlp}s for letters and two recurrent models for words
(one \gls{bilstm} with custom temporal attention for \gls{asl},
one stacked Bi-/Uni-\gls{lstm} with TimeDistributed spatial encoder for
\gls{arsl}). The models were exposed through a FastAPI backend with
both \gls{rest} and \gls{ws} endpoints, fronted by a React web client and
a React-Native (Expo) mobile client; user accounts and chat are handled
by a separate Node.js service. The system is containerised with Docker
and deployed to Railway (backend) and Vercel (web).

\section{Contributions to the Field}
\label{sec:conc:contrib}

The four contributions stated in \cref{sec:intro:contributions} were
realised and quantitatively validated:
\begin{itemize}
  \item \textbf{C-1} The unified bilingual pipeline is delivered as
        an open-source repository.
  \item \textbf{C-2} The ArSL Word v2 architecture outperforms a
        hands-only \gls{bilstm} baseline on \gls{karsl}-502
        (see \cref{tab:ablate}).
  \item \textbf{C-3} The custom \texttt{TemporalAttention} layer
        provides a measurable Top-5 gain at $<\,1$\,\% parameter cost
        (\cref{tab:ablate}).
  \item \textbf{C-4} The mobile client runs \gls{tflite} models on-device
        at $\geq 25$~\gls{fps} with a model footprint $\leq 5$~MB each
        (\cref{tab:rt}).
\end{itemize}

\section{Practical Impact}
\label{sec:conc:impact}

The most direct beneficiaries are deaf and hard-of-hearing Arabic
speakers, who gain a free, open, and bidirectional communication aid.
The system has potential application in special-needs education,
telemedicine, public-service kiosks, and customer service. As an
open-source release, it also lowers the entry barrier for future
\gls{arsl} research.
